\CQUsection{绪论}

\CQUsubsection{研究背景与意义}
脑血管疾病长期位居我国居民死因构成前列，其中颅内动脉瘤即 Intracranial Aneurysm, IA 可在无症状阶段被影像学检查偶然发现；一旦瘤壁完整性受损并导致蛛网膜下腔出血，病情通常较为凶险，遗留神经功能障碍的比例也不容忽视。数字减影血管造影即 DSA 被视作诊断“金标准”，但它的有创性以及费用限制了其在大规模筛查场景中的普及。CT 血管成像，即 CT Angiography, CTA 具有无创、可及性高的特性，在当前门诊以及急诊情境下，是动脉瘤检出与术前评估的重要工具\cite{chen2021,ou2022}。

然而在临床实际工作当中，仍然存在三方面痛点。第一，单次检查产生的切片可达数百层，医生需要在三维空间中追踪复杂血管脉络，并且辨别细微隆起，长时间阅片容易引起视觉疲劳以及漏诊。第二，不同扫描协议、重建层厚以及对比剂充盈程度会带来图像对比度、颜色以及间隔等差异，同一算法在不同数据以及硬件条件下难以保持足够稳定性与可靠泛化能力。第三，是否建议干预不仅取决于“有没有瘤”，还与瘤体大小、形态不规则度、位置、患者年龄以及合并症等因素相关；现有临床诊断往往还需要参考 PHASES、ELAPSS 以及日本 UCAS 等研究所揭示的自然病程与风险规律\cite{greving2014,backes2017,morita2012,thompson2015}。因此，如果仅依靠分割掩膜而缺少风险预测，实际上较难直接支撑完整的临床规划任务。

近年来，以多层神经网络为主的深度学习在医学影像分割任务中取得了显著进展。全卷积网络以及 U-Net 族结构能够输出像素级病灶概率图，极大程度上缩短临床医师的标注辅助时间\cite{unet,long2015}；三维卷积网络可直接建模沿血管走向的上下文信息，同时有利于识别瘤体在层间连续性等形态学特征\cite{unet3d,milletari2016}。Attention U-Net 在解码路径中加入注意力门控单元，使网络更关注与任务相关的空间区域，从而提升模型表现，且已在腹部脏器、胰腺等分割任务中得到验证\cite{attention_unet}。nnU-Net 的成功进一步说明，在医学分割任务中，除网络骨架外，数据预处理、补丁采样以及后处理同样会对性能产生重要影响\cite{nnunet}。在脑血管领域，已有工作尝试把粗定位与细分割结合起来使用，或借助多尺度三维卷积来提升细小病灶检出率\cite{kamnitsas2017,chen2021,ou2022}。与此同时，运用动脉瘤临床信息以及动脉瘤本体属性来预测风险的网络也开始出现，这提示出“影像 + 临床”联合建模的重要性\cite{strokenet,xiong2025}。在表格数据学习领域，TabTransformer 与 FT-Transformer 等模型证明了注意力机制不仅适用于视觉特征建模，同时也适用于多维临床字段之间的交互建模，这为本项目风险预测模块的设计提供了参考\cite{tabtransformer2020,fttransformer2021}。

在这一背景下，本文依托课程设计任务，围绕“动脉瘤的自动分割定位”以及“破裂风险的概率预测”两个相互衔接的子问题，完成一套以配置文件驱动、脚本可复现为特征的软件系统。该系统一方面可以减轻医师在体数据层面的重复勾画负担，另一方面也为后续把影像组学特征或形态学参数与表型数据进行融合提供工程接口，具有较明确的辅助诊疗价值，并为医学影像与临床信息结合的医学软件系统提供完整解决方案。

\CQUsubsection{国内外研究现状}
\CQUsubsubsection{动脉瘤与脑血管影像分析的传统途径}
在深度学习普及之前，血管分割以及病灶提取主要还是依靠传统图像处理手段，比如基于强度的阈值分割、区域生长、水平集方法，或者部分基于几何先验的滤波器。此类流程往往需要人工去调节种子点、平滑核大小或曲率约束参数，并且对低对比度瘤颈、邻近骨质高密度灶等干扰较为敏感，在面对不同数据集时常常需要重新调参。面向风险评估，临床诊断多采用 Hunt-Hess 分级、WFNS 评分、PHASES 评分或 ELAPSS 评分等规则体系，其中 PHASES 更侧重未破裂动脉瘤未来破裂风险，ELAPSS 更强调肿瘤增长风险预测\cite{greving2014,backes2017}。这些评分简洁且具有可解释性，便于临床沟通，但对瘤体形态不规则、局部血管形态变化以及影像纹理等非线性特征的捕捉与运用能力较为有限。统计学习阶段则常选用逻辑回归、随机森林等传统机器学习模型；它们在变量个数较少、线性可分性较强的回顾性队列上可以给出可解释结果，但对多维关联特征的自动建模能力仍然有限。

\CQUsubsubsection{基于深度学习的分割与多模态预测}
随着公开数据集以及算力条件的改善，U-Net 及其三维变体逐渐成为血管与病灶分割的主流骨架\cite{unet,unet3d}。nnU-Net 等工作进一步强调预处理、网络拓扑以及后处理的自动化配置，从而提升跨任务迁移效率\cite{nnunet}。针对颅内动脉瘤，研究者探索了从血管分割到瘤体检测的多阶段流程，或引入注意力以及多尺度模块来强调细小结构\cite{chen2021,ou2022}。在预测层面，Transformer 与注意力机制为序列化或多模态特征融合提供了较标准的实现方式\cite{vaswani2017}；StrokeNet 等多任务结构尝试同时编码影像以及临床信息，以预测卒中结局\cite{strokenet}。这些工作表明，分割精度与下游临床推理并不是割裂的——高质量掩膜既可以直接服务于手术规划，同时也可以作为形态学等衍生特征的载体。
\CQUsubsection{本文面临的问题与主要思路}
综合文献以及工程经验，本文把任务难点概括为三点。（1）\textbf{体数据占用与显存矛盾}：完整 CTA 体积可达数百兆字节以上，如果一次性加载多病例批次，普通 GPU 往往难以承受；因此需要在“流式 IO—补丁采样—梯度更新”之间取得平衡。（2）\textbf{前景极度稀疏}：动脉瘤体素在整幅体积中的占比极低，单纯交叉熵容易使模型偏向预测背景；因此需要区域重叠型损失以及采样策略共同进行约束。（3）\textbf{风险标签与特征异质性}：破裂与未破裂病例分布不均衡，类别字段存在长尾取值；因此需要合理的加权损失、阈值搜索以及特征筛选，避免模型仅学习到训练集分布而导致过拟合。

针对上述难点，本文采用“分割骨干 + 表格风险头”的系统级方案：分割侧选用三维 Attention U-Net，并且配合病例级流式训练与前景优先补丁采样；风险侧使用基于交叉注意力的表格模型，统一处理数值以及类别特征，并在验证集上搜索决策阈值。

\CQUsubsection{课题来源与研究内容}
本课题来源于指导教师下达的毕业设计任务书，目标是对颅内动脉瘤 CTA 数据开展计算机辅助检测，并在临床表型与影像衍生特征基础上完成破裂风险的二分类预测。本文工作具体包括：
\begin{enumerate}
	\item 基于 \texttt{AttentionUnet} 实现三维分割网络，完成编码器—解码器通道配置、注意力门控以及 sigmoid 输出的端到端前向传播；
	\item 在 \texttt{MedicalPatchDataset} 上实现多目录数据扫描、重采样、标准化以及前景优先/顺序采样策略，支撑低显存条件下的长时间训练；
	\item 基于 Dice/Tversky/Focal 与 BCE 的组合损失、Adam/AdamW/SGD 以及余弦退火等配置化优化流程，训练过程写入 TensorBoard，并按间隔保存最新模型与验证集最优模型；
	\item 推理阶段把补丁结果拼接为完整体积，恢复 spacing、origin、direction，输出概率体与二值掩膜，供可视化或形态学特征提取使用；
	\item 构建表格风险数据的清洗、互信息特征筛选、加权 BCE 训练、验证集 F1 最优阈值搜索以及测试集报告导出流程；
	\item 提供 FastAPI 交互界面，用于文件上传、分割以及风险预测流程。
\end{enumerate}

\CQUsubsection{论文结构安排}
全文共分五章。第一章介绍研究背景、国内外现状、问题难点以及研究内容。第二章梳理医学图像分割、三维卷积网络、注意力机制、类别不平衡学习以及表格建模等相关理论基础。第三章从系统架构出发，分别说明分割模块、风险模块的模型设计与实现要点，并简要介绍 Web 服务接口。第四章介绍实验环境、数据集与预处理约定、评价指标以及结果分析框架，说明如何依据日志与配置文件复现指标。第五章总结全文，并展望未来改进方向。

\CQUsubsection{本章小结}
本章从临床需求以及技术趋势两个方面阐述了颅内动脉瘤辅助分析的研究动机，综述了传统方法与深度学习方法的优劣，归纳出本文面向的三个核心难点及对应思路，并列出与源代码模块相对应的研究内容。后续章节将进一步展开算法与系统实现细节。
